\section{The Base Anatomy}

This appendix is the complete specification of the base. It gives the full knit, every collision with its complete transition table, the exact layout of the state buffer, and the geometry of the direction field. Everything here is discrete. No real number and no decimal appears anywhere in the base. Continuity, in the form of sound, the Lorentz symmetry, smooth gravity, and the Lie groups, is only ever emergent. It is never fundamental.

The universe is a vibe mesh, a flow of eight atoms. The single substance is the vibe, the experience carried in each site and valued by a tone. The whole, the eight atoms running as one, is the flow. The aim here is to write down those eight, the law that turns them, and nothing else.

\subsection{The axioms}

The axioms are constraints on what the base is allowed to contain. They are load-bearing. Breaking one is a cheat, and the model is judged in part by its refusal to cheat. They are stated first so that every later choice can be checked against them.

\begin{axiom}[The eight base axioms]
The base satisfies, all at once, the eight constraints of Table~\ref{tab:axioms}. They fix it as discrete, deterministic, reversible in the bulk, charge-conserving, local, free of hidden state, monistic, and separated from physics by an emergence gap.
\end{axiom}

\begin{table}[ht]
\centering
\begin{tabular}{@{}lll@{}}
\toprule
 & axiom & meaning \\
\midrule
A1 & discrete & no real numbers anywhere, finite mesh per beat, finite directions, ternary values, discrete time, finite symmetry \\
A2 & deterministic & the model and every initial condition are fixed functions, never a dice roll, we vary size not seeds \\
A3 & reversible bulk & collide then stream is an invertible permutation of states, time can run backward in the closed bulk \\
A4 & charge-conserving & the knit conserves the total tone, the sum of all slots \\
A5 & local & a dock's update reads only its own slots, and streaming moves one dock \\
A6 & no hidden state & a dock is exactly its directional tones, with no labels, flags, or memory \\
A7 & monism & the only entity is the vibe, one substance with two motions, bound is matter and waving is light \\
A8 & emergence gap & raw tones are not physics, physics appears only after one to three emergent coarse-graining layers \\
\bottomrule
\end{tabular}
\caption{The eight axioms that constrain the base.}
\label{tab:axioms}
\end{table}

The deepest consequence of A1 together with A8 is the two-layer reading of the model. The base is discrete, and we call it Layer~$0$. The self and every recognizable piece of physics are emergent, and live at Layer~$1$ and above, as the coarse-grained average of many discrete grains.

\begin{principle}[Two layers]
The base is discrete at Layer~$0$. Physics, including the self, is the coarse-grained average of many discrete grains at Layer~$1$ and above. Nothing at Layer~$0$ is read directly as a particle, a force, or a smooth field.
\end{principle}

\subsection{The eight atoms}

The base is exactly eight atoms. Each is a soft four-letter word starting with a distinct letter, so each is also a clean single-letter symbol for the mathematics. Between them they answer a small fixed set of questions. Where is anything, when does anything happen, what is the substance, how does the state change, and how does the space change. Adding a ninth atom as a base thing is a cheat.

\begin{definition}[The eight atoms]
The base is the eight atoms mesh, dock, site, tone, lean, knit, wake, and beat, with the single-letter symbols $m$, $d$, $s$, $t$, $l$, $k$, $w$, $b$. The whole they assemble into is the flow, written $F$.
\end{definition}

\begin{table}[ht]
\centering
\begin{tabular}{@{}llll@{}}
\toprule
atom & letter & plain meaning & the commitment it carries \\
\midrule
mesh & $m$ & the whole space & the $\{3,4,3,4\}$ hyperbolic honeycomb, a weave of docks, finite at any beat and ever-growing \\
dock & $d$ & a here-now & the geometric cell, a place of transit, each beat tones arrive, knit, and move on \\
site & $s$ & one direction out of a dock & one of the $24$ $D_4$ directions, the directed slot where a tone lives \\
tone & $t$ & the substance value & one ternary value in $\{-1, 0, +1\}$ on each of a dock's $24$ sites \\
lean & $l$ & the value order & the order $-1 < 0 < +1$, the tilt that gives charge and the create move its direction \\
knit & $k$ & the state-law & a reversible, charge-conserving local collision, including the create move, then streaming \\
wake & $w$ & the space-law & new docks born at peace at the open, ever-expanding edge \\
beat & $b$ & the discrete step & discrete, synchronous, one beat is collide then stream, invertible, no built-in direction \\
\bottomrule
\end{tabular}
\caption{The eight atoms of the base.}
\label{tab:atoms}
\end{table}

The eight group into a small natural furniture. There is a where, carried by mesh, dock, and site. There is a what, the substance, carried by tone and lean. There is a how the state changes, carried by knit. There is a how the space changes, carried by wake. And there is a when, carried by beat. Two laws appear rather than one, because two distinct things evolve. The state evolves under the knit, and the space itself evolves under the wake.

Two terms that a reader might expect to find here are absent on purpose. The arrow, the direction of time, is not an atom. It emerges, the wake acting through the lean and the create move, and is taken up in the subsection on the wake below. The attraction, the binding by radiation, is also not an atom. It emerges, and is taken up in the final subsection.

\subsection{The mesh and its two regions}

The base space is the $\{3,4,3,4\}$ hyperbolic honeycomb, a regular tiling of four-dimensional hyperbolic space. Its geometry is fixed, not dynamical. The mesh does not bend in response to matter. The bending that general relativity sees is emergent, taken up later in the paper. The extent of the mesh is finite at any beat and ever-growing. At any moment it is a finite region with a real boundary, the wake, which moves outward as the mesh grows.

\begin{table}[ht]
\centering
\begin{tabular}{@{}lll@{}}
\toprule
region & what it is & role \\
\midrule
bulk & the negatively-curved hyperbolic interior & a stronger bath, it disperses structure exponentially, it does not bind \\
cusp & a parabolic point at infinity, its cross-section a horosphere & where physics lives, the horosphere is a flat Euclidean lattice \\
\bottomrule
\end{tabular}
\caption{The two regions of the honeycomb.}
\label{tab:regions}
\end{table}

Physics lives at the cusp, on a horosphere whose cross-section is the flat $D_4$ lattice, a four-dimensional integer lattice. The perfectly flat infinite cusp is a limit, not a state the mesh is ever in at a finite beat. At any finite beat the mesh is a finite horosphere patch at finite cusp-height, slightly curved, with a real wake. Eternal expansion is the patch flowing deeper toward the cusp, growing and flattening each beat, forever approaching the ideal flat $D_4$ space but never reaching it.

\begin{principle}[Finite and growing, never finished]
At every finite beat the mesh is a finite, slightly curved horosphere patch with a real wake. The perfectly flat, infinite $D_4$ cusp is the attractor it forever approaches and never reaches. Both statements hold, at different times.
\end{principle}

\subsection{The sites, the 24 directions}

Each dock has $24$ neighbours, reached along the $24$ $D_4$ root vectors, every coordinate arrangement of $(\pm 1, \pm 1, 0, 0)$ in four dimensions. This $24$-direction set is the $24$ sites. Every direction has the same length, with norm squared $2$, so the sites are single-speed. There is no slow direction and no fast direction in the base.

The full set of $24$ directions, grouped by which two of the four axes $x, y, z, w$ are nonzero, is laid out in Table~\ref{tab:directions}.

\begin{table}[ht]
\centering
\begin{tabular}{@{}lll@{}}
\toprule
index & nonzero axes & the four root vectors \\
\midrule
$0,1,2,3$ & $(x,y)$ & $(+,+,0,0)\ (+,-,0,0)\ (-,+,0,0)\ (-,-,0,0)$ \\
$4,5,6,7$ & $(x,z)$ & $(+,0,+,0)\ (+,0,-,0)\ (-,0,+,0)\ (-,0,-,0)$ \\
$8,9,10,11$ & $(x,w)$ & $(+,0,0,+)\ (+,0,0,-)\ (-,0,0,+)\ (-,0,0,-)$ \\
$12,13,14,15$ & $(y,z)$ & $(0,+,+,0)\ (0,+,-,0)\ (0,-,+,0)\ (0,-,-,0)$ \\
$16,17,18,19$ & $(y,w)$ & $(0,+,0,+)\ (0,+,0,-)\ (0,-,0,+)\ (0,-,0,-)$ \\
$20,21,22,23$ & $(z,w)$ & $(0,0,+,+)\ (0,0,+,-)\ (0,0,-,+)\ (0,0,-,-)$ \\
\bottomrule
\end{tabular}
\caption{The $24$ $D_4$ root directions, the sites of a dock.}
\label{tab:directions}
\end{table}

The count is $6$ axis-pairs times $4$ sign-choices, which is $24$. These are exactly the $D_4$ roots, the vertices of the $24$-cell, and the unit Hurwitz quaternions, the same $24$ objects in three guises. The quaternion reading is the one that matters for the self.

The symmetry group of the sites is $F_4$, of order $1152$, the symmetry group of the $24$-cell. It is finite. This high symmetry is what lets continuous rotational symmetry be restored in the emergent infrared, the discrete $F_4$ flowing to a continuous group only as a coarse-grained limit, never in the base.

\begin{proposition}[The site set is the $D_4$ root system]
The $24$ directions out of a dock are the $D_4$ root vectors, equivalently the vertices of the $24$-cell and the unit Hurwitz quaternions. Their symmetry group is the finite group $F_4$ of order $1152$. All $24$ are single-speed, with norm squared $2$.
\end{proposition}

\subsubsection{The opposite map and the lines}

Streaming needs a well-defined reverse direction. The opposite of a direction is its negation. Every direction pairs with exactly one opposite, so the $24$ directions form $12$ opposite-pairs, which we call lines. The two slots of a line carry the two tones moving head-on through the dock. A line is the unit the collision acts on.

Each axis-pair block of four directions splits into two lines, $(+,+)$ with $(-,-)$ and $(+,-)$ with $(-,+)$, giving $6$ blocks times $2$ lines, which is $12$ lines. For example, $(+,+,0,0)$ is opposite $(-,-,0,0)$, and the two are one line. Distance on the mesh is the number of single-direction hops, which is the natural radius for a structure.

\subsection{The tone and the state buffer}

The only substance is the tone, a ternary value on each directed slot. It is the vibe experience itself, valued. A single tone is one quale, read as pain, peace, or pleasure. A whole dock, its $24$ tones together, is the hyper-color, the richer local feel built from the $24$ single qualia.

\begin{table}[ht]
\centering
\begin{tabular}{@{}llll@{}}
\toprule
value & name & feeling & meaning \\
\midrule
$-1$ & pain & aversion & a unit charge moving along that direction, negative sign \\
$0$ & peace & neutral & empty, no charge in that slot \\
$+1$ & pleasure & attraction & a unit charge moving along that direction, positive sign \\
\bottomrule
\end{tabular}
\caption{The three tones, the single quale per site.}
\label{tab:tones}
\end{table}

The value ordering $-1 < 0 < +1$ is the lean, an atom in its own right. It is what makes peace the middle, the empty state from which a balanced pair can be born, and it is the order through which the emergent arrow runs.

\begin{result}[A dock holds exactly $24$ tones]
A dock holds exactly $24$ tones, one per site. There is no rest slot, no twenty-fifth direction, and no home slot. Binding is supplied later by an emergent attraction, not by a still center in the base.
\end{result}

The whole world is one flat array of ternary values, and that array is the universe. Nothing else is stored. The only allowed values are minus one, zero, and plus one. Its length is the number of docks times the degree, where the degree is twenty-four, the number of sites. A dock's full state is its twenty-four ternary values, one per site, so each dock is twenty-four trits.

\begin{table}[ht]
\centering
\begin{tabular}{@{}ll@{}}
\toprule
object & definition \\
\midrule
dock & one place with twenty-four directed slots, one tone each \\
slot & the tone in one direction at one dock \\
line & a direction and its opposite, the slot pair, the two tones moving head-on \\
dock tone & the sum of a dock's slots, its net local charge \\
charge & the sum of every slot over the whole mesh, conserved by the knit \\
\bottomrule
\end{tabular}
\caption{The objects of the state buffer.}
\label{tab:buffer}
\end{table}

\subsubsection{Matter and radiation, two motions of one tone}

There is one tone, not two. Matter and light are two motions of it. This is mass-energy equivalence built into the base. A charge that is bound and localized is matter, a trapped persistent lump of tone, which is mass as trapped energy. The same tone, spread as a net-zero long-wavelength ripple of many tones, is light, a propagating collective disturbance, a photon read as a phonon. There is no separate matter field and no separate radiation field.

\begin{table}[ht]
\centering
\begin{tabular}{@{}lll@{}}
\toprule
motion & what it is & physics name \\
\midrule
a charge bound and localized & a trapped, persistent lump of tone & matter, mass, trapped energy \\
a net-zero, long-wavelength ripple of many tones & a propagating collective disturbance & light, sound, a photon as a phonon \\
\bottomrule
\end{tabular}
\caption{One tone, two motions.}
\label{tab:motions}
\end{table}

\subsection{The collision, every table}

A collision is a local, in-place map on one dock's directional slots. It must satisfy four requirements, set out in Table~\ref{tab:requirements}. It reads and writes only the one dock, it is a reversible permutation of the dock's states, it conserves the per-line tone sum, and it acts on each line independently.

\begin{table}[ht]
\centering
\begin{tabular}{@{}ll@{}}
\toprule
requirement & meaning \\
\midrule
local & reads and writes only the one dock's slots \\
reversible & a permutation of the dock's states, invertible \\
charge-conserving & the per-line tone sum is unchanged \\
acts per line & the committed knit applies a $9$-state table to each opposite-pair independently \\
\bottomrule
\end{tabular}
\caption{What a collision must be.}
\label{tab:requirements}
\end{table}

A line carries two tones, so it has $3 \times 3 = 9$ states. A per-line collision is a permutation of those nine states. We write a line as $(\text{left}, \text{right})$, the tones in a direction and its opposite.

\subsubsection{The committed base collision}

This is the knit. It is the create-flip-annihilate cycle plus a hop plus two inert fixed points. It runs on every one of the $12$ lines independently. The forward table is Table~\ref{tab:forward}.

\begin{table}[ht]
\centering
\begin{tabular}{@{}lllll@{}}
\toprule
 & in $(\text{left}, \text{right})$ & out $(\text{left}, \text{right})$ & move-type & plain meaning \\
\midrule
$1$ & $(-1,-1)$ & $(-1,-1)$ & inert & like charges coexist, nothing happens \\
$2$ & $(+1,+1)$ & $(+1,+1)$ & inert & like charges coexist, nothing happens \\
$3$ & $(-1,0)$ & $(0,-1)$ & hop & the lone $-1$ charge moves across the dock \\
$4$ & $(0,-1)$ & $(-1,0)$ & hop & the lone $-1$ charge moves across the dock \\
$5$ & $(+1,0)$ & $(0,+1)$ & hop & the lone $+1$ charge moves across the dock \\
$6$ & $(0,+1)$ & $(+1,0)$ & hop & the lone $+1$ charge moves across the dock \\
$7$ & $(0,0)$ & $(+1,-1)$ & create & the lean, peace brings out a balanced pair \\
$8$ & $(+1,-1)$ & $(-1,+1)$ & flip & the balanced pair reverses \\
$9$ & $(-1,+1)$ & $(0,0)$ & annihilate & the reversed pair returns to peace \\
\bottomrule
\end{tabular}
\caption{The forward knit, every one of the nine line states.}
\label{tab:forward}
\end{table}

The structure is two inert fixed points, rows $1$ and $2$, two hop swaps that transport a lone charge, rows $3$ to $6$, and one three-cycle $(0,0) \to (+1,-1) \to (-1,+1) \to (0,0)$ that is the create move along the lean, rows $7$ to $9$. It is a permutation of the nine states, so it is reversible, and every row preserves $\text{left} + \text{right}$, so it is charge-conserving.

\begin{theorem}[The knit is a reversible charge-conserving permutation]
The forward table of the knit is a permutation of the nine line states. The hops and inert states are self-inverse. The three-cycle inverts to $(0,0) \to (-1,+1) \to (+1,-1) \to (0,0)$. Every row preserves the line sum $\text{left} + \text{right}$. Hence the knit is reversible and charge-conserving, satisfying A3 and A4 line by line.
\end{theorem}

To run time backward, the knit is replaced by its inverse, Table~\ref{tab:inverse}. The hops and inert states are unchanged, since they are self-inverse. Only the three-cycle reverses.

\begin{table}[ht]
\centering
\begin{tabular}{@{}ll@{}}
\toprule
in $(\text{left}, \text{right})$ & out $(\text{left}, \text{right})$ \\
\midrule
$(-1,-1)$ & $(-1,-1)$ \\
$(+1,+1)$ & $(+1,+1)$ \\
$(-1,0)$ & $(0,-1)$ \\
$(0,-1)$ & $(-1,0)$ \\
$(+1,0)$ & $(0,+1)$ \\
$(0,+1)$ & $(+1,0)$ \\
$(0,0)$ & $(-1,+1)$ \\
$(+1,-1)$ & $(0,0)$ \\
$(-1,+1)$ & $(+1,-1)$ \\
\bottomrule
\end{tabular}
\caption{The inverse knit, used to run time backward.}
\label{tab:inverse}
\end{table}

The behavior of the committed knit is summarized in Table~\ref{tab:behavior}. It is reversible and charge-conserving. It does not conserve momentum, because it reflects, and that reflection is exactly why it pins a charge and seals the radiation channel. It carries the create move, the active vacuum of the lean. It is good for the create move and for a confined identity, and weak for a self that must radiate, since the pinning seals the channel a self would shed its disturbance through.

\begin{table}[ht]
\centering
\begin{tabular}{@{}ll@{}}
\toprule
property & value \\
\midrule
reversible & yes \\
conserves charge & yes \\
conserves momentum & no, it reflects, which is why it pins and seals radiation \\
has the create move & yes, the lean's create, an active vacuum \\
good for & the create move, a confined identity \\
weak for & a self that must radiate, the pinning seals the radiation channel \\
\bottomrule
\end{tabular}
\caption{Behavior of the committed knit.}
\label{tab:behavior}
\end{table}

\subsubsection{The variant collisions}

Each variant is a reversible, charge-conserving collision built to test one question. None is the committed base. They are catalogued in full so the machinery is documented and the one open choice is clear.

The first is pass-through, the trivial control. It is the identity, leaving every input unchanged, so every slot streams forever with no interaction.

The second is momentumRotate2D, the square reference. On the four-direction square set of sites, in order east, west, north, south, a zero-momentum head-on pair on one axis rotates to the other. If $(E,W,N,S)$ is $(s,s,0,0)$ with $s \neq 0$ it becomes $(0,0,s,s)$, and if it is $(0,0,s,s)$ it becomes $(s,s,0,0)$. Anything else, a lone charge or a config with net momentum, streams unchanged. It is a reversible involution that conserves charge and momentum, the canonical two-dimensional reference collision.

The third is headOnRotate, the mobility probe. It is the generalization of the square reference to any set of sites. It pairs the $12$ lines two at a time into $6$ line-pairs, and moves a zero-momentum head-on pair $(s,s)$ from a full line to its empty partner line, as in Table~\ref{tab:headon}. A lone particle, or any config with net momentum, is left untouched, so it streams ballistically. This is exactly the mobility the pinning pair table denies. It is a reversible involution that conserves charge and momentum, but it has no create move. This is the collision the self work favors, with the create move dropped and binding supplied instead by the emergent attraction.

\begin{table}[ht]
\centering
\begin{tabular}{@{}ll@{}}
\toprule
condition for a line-pair $A$, $B$ & action \\
\midrule
$A$ is $(s,s)$, $s \neq 0$, and $B$ is $(0,0)$ & $A$ becomes $(0,0)$, $B$ becomes $(s,s)$ \\
$B$ is $(s,s)$, $s \neq 0$, and $A$ is $(0,0)$ & $B$ becomes $(0,0)$, $A$ becomes $(s,s)$ \\
otherwise & both lines unchanged \\
\bottomrule
\end{tabular}
\caption{headOnRotate, the mobility probe.}
\label{tab:headon}
\end{table}

The fourth is bindAndMove, the pair table with the hop turned off. The create-flip-annihilate cycle is kept, but the lone-charge hop is off, so lone charges stream, as in Table~\ref{tab:bindmove}. It is reversible and charge-conserving. It is mobile, but the create move destabilizes the vacuum globally, so it gives no clean moving bound state.

\begin{table}[ht]
\centering
\begin{tabular}{@{}llll@{}}
\toprule
 & in $(\text{left}, \text{right})$ & out $(\text{left}, \text{right})$ & move-type \\
\midrule
$1$ & $(-1,-1)$ & $(-1,-1)$ & inert \\
$2$ & $(+1,+1)$ & $(+1,+1)$ & inert \\
$3$ & $(-1,0)$ & $(-1,0)$ & hop off, the lone charge streams \\
$4$ & $(0,-1)$ & $(0,-1)$ & hop off \\
$5$ & $(+1,0)$ & $(+1,0)$ & hop off \\
$6$ & $(0,+1)$ & $(0,+1)$ & hop off \\
$7$ & $(0,0)$ & $(+1,-1)$ & create \\
$8$ & $(+1,-1)$ & $(-1,+1)$ & flip \\
$9$ & $(-1,+1)$ & $(0,0)$ & annihilate \\
\bottomrule
\end{tabular}
\caption{bindAndMove, the pair table with the hop off.}
\label{tab:bindmove}
\end{table}

The fifth is leakyConfine, matter and photon in one field. It is a single move that confines charged line-states, so a body keeps its identity, while letting charge-neutral line-states stream, so neutral ripples, the photon, propagate, as in Table~\ref{tab:leaky}. There is no create move, so the vacuum stays empty and there is no churn. It is a self-inverse involution, reversible and charge-conserving. The matter is the tone bound, the photon is the tone waving, in one field.

\begin{table}[ht]
\centering
\begin{tabular}{@{}llll@{}}
\toprule
 & in $(\text{left}, \text{right})$ & out $(\text{left}, \text{right})$ & move-type \\
\midrule
$1$ & $(-1,-1)$ & $(-1,-1)$ & inert \\
$2$ & $(+1,+1)$ & $(+1,+1)$ & inert \\
$3$ & $(+1,0)$ & $(0,+1)$ & reflect, the lone charge bounces inward, confine \\
$4$ & $(0,+1)$ & $(+1,0)$ & reflect \\
$5$ & $(-1,0)$ & $(0,-1)$ & reflect \\
$6$ & $(0,-1)$ & $(-1,0)$ & reflect \\
$7$ & $(0,0)$ & $(0,0)$ & peace left alone, no create move \\
$8$ & $(+1,-1)$ & $(+1,-1)$ & stream, a neutral pair propagates, the phonon \\
$9$ & $(-1,+1)$ & $(-1,+1)$ & stream \\
\bottomrule
\end{tabular}
\caption{leakyConfine, matter and photon in one field.}
\label{tab:leaky}
\end{table}

The sixth is stickyReflect, the capture probe. It is a crowding-dependent reflection that reads the whole dock, not a single line. If fewer than two slots are nonzero, a lone particle or empty dock, the dock is left unchanged and streams. If two or more slots are nonzero, a collision, it reflects every line, swapping each slot with its opposite. It is a reversible involution that conserves charge. The reflection of a head-on collision is elastic, the charges exchange and reverse and then leave, so it does not capture.

\subsubsection{The collision summary}

Table~\ref{tab:collisions} gathers all the collisions side by side. The pair table is the committed base. Every other entry is a control or a probe.

\begin{table}[ht]
\centering
\begin{tabular}{@{}llllll@{}}
\toprule
collision & label & reversible & charge & momentum & key property \\
\midrule
pass-through & control & yes & yes & yes & no interaction \\
\textbf{pair table} & \textbf{base} & yes & yes & no & the create move (lean), pins and confines \\
momentumRotate2D & reference & yes & yes & yes & 2D head-on rotate \\
headOnRotate & variant & yes & yes & yes & mobile, elastic scatter, no create move \\
bindAndMove & variant & yes & yes & no & create kept, hop off, vacuum-unstable \\
leakyConfine & variant & yes & yes & no & confine charge, stream neutral \\
stickyReflect & variant & yes & yes & no & crowd-reflect, elastic, no capture \\
\bottomrule
\end{tabular}
\caption{Every collision, the committed base and the probes.}
\label{tab:collisions}
\end{table}

\begin{result}[The one open choice]
The single open choice the model carries is which collision the self uses. It is the pinning pair table, which carries the create move but seals radiation, versus the radiating momentum-conserving collision, which has no create move. Both are reversible and charge-conserving. The self work favors the momentum-conserving collision with binding supplied by the emergent attraction.
\end{result}

\subsection{Streaming and the beat}

After the collision, each tone moves one dock along its own direction. Streaming touches no values, it only moves them. It is a bijection, so it is exactly reversible. Read the forward stream as follows. The tone now in a direction at this dock came from the neighbour on the opposite side, still traveling that way. To run streaming backward, take the tone in a direction from the neighbour on that same side instead. A tone keeps its line of travel, which is why the opposite map must be well defined.

One beat is the discrete unit of time. It is two steps, applied synchronously to every dock. First, collide, applying the collision to every dock's slots in place. Second, stream, moving each slot one dock along its direction. There is no continuous time and no real-valued clock. Time is the integer count of beats.

Because both halves are invertible, the whole beat is invertible. To run one beat backward, un-stream then un-collide, using the paired inverse table for a non-involutive collision like the pair table, or the same table for an involution. So the closed bulk has no arrow of time. Forward and backward are equally valid. The arrow lives entirely in the wake and the open boundary, never in the beat itself.

\begin{principle}[The complete base dynamics in four moves]
Every beat, in order, the base does four things. First it collides, replacing each line's two tones by the knit's output. Second it streams, moving every tone one dock along its direction. Third it grows, new docks waking into being at the open edge, each born at peace. Fourth it bathes, whatever streamed off the open edge is gone for good and never returns. Everything else in physics is what this grows into.
\end{principle}

\subsection{The wake, the create move, and the emergent arrow}

The arrow, the direction of time, is not a base atom. It emerges from three pieces that already live among the eight atoms. The lean supplies the value order $-1 < 0 < +1$, which defines charge and gives the create move a direction. The create move, part of the knit, lets peace bring out a balanced pair, which then flips and annihilates in a reversible three-cycle. This fills empty space with fluctuating pairs, an active vacuum, whose only force at a distance is a weak depletion attraction, too weak to bind on its own. The wake births new docks at peace at the open edge, so the mesh grows and the open region stays empty out there. Radiation that leaves never returns. This is the bath.

A closed, toroidal mesh cannot host a self. It recurs, by the Poincare recurrence, and the bare knit on a closed system produces only churn. The wake is the low-entropy source the arrow needs, and it is not optional.

\begin{result}[The arrow is emergent, the beat is base]
The beat ticks with no built-in direction, since the bulk is reversible. The arrow is the thermodynamic direction produced by the lean, the create move, and the wake acting together. The wake is the indispensable low-entropy source. The beat is base, the arrow is emergent.
\end{result}

\subsection{The attraction, emergent not base}

The eight atoms give a self three things. They give it identity, a conserved charge. They give it agency, the ability to radiate a disturbance to the bath and to heal. And they give it a bound body, the radiation-pressure shadow of a vacuum-excluding mass. The attraction that binds the body does not enter the base.

The binding is the radiation-pressure shadow of the active vacuum. The create move, part of the knit, fills empty space with fluctuating pairs. A mass excludes the vacuum, so the pressure on it is unbalanced and presses two bodies together. The wake carries the radiation away one-way, so the binding is permanent. This is built entirely from the knit's create move plus the open growing mesh and its wake, so it is emergent, not a base ingredient. An experiment confirmed that this shadow pressure confines a single-speed body. The self binds as a cluster of moving tones. Its identity is the topological winding read at $24$-direction resolution, and its binding is the shadow pressure. A dock has $24$ sites and needs no rest slot.

\begin{result}[The attraction is emergent]
The attraction is the radiation-pressure shadow of the active vacuum on the open growing mesh. It is built from the knit's create move, the wake, and the exclusion of vacuum by mass, all of which are already in the base. It is therefore emergent, not a ninth atom. No rest slot and no new field are added to the base.
\end{result}

The ledger of Table~\ref{tab:ledger} closes the appendix. The eight atoms are base. The arrow and the attraction are emergent, produced by those eight and never grafted on.

\begin{table}[ht]
\centering
\begin{tabular}{@{}llll@{}}
\toprule
atom & category & what it is & status \\
\midrule
mesh & where & $\{3,4,3,4\}$ honeycomb, weave of docks, cusp a flat $D_4$ horosphere, open and growing & base atom \\
dock & where & the here-now cell, a place of transit, $24$ sites & base atom \\
site & where & one of the $24$ $D_4$ directions out of a dock, the directed slot & base atom \\
tone & what & ternary $\{-1,0,+1\}$ on each of $24$ sites per dock & base atom \\
lean & what & the value order $-1 < 0 < +1$, charge and the create's direction & base atom \\
knit & state-law & reversible charge-conserving collide per-line, including create, then stream & base atom \\
wake & space-law & new docks born at peace at the open, ever-expanding edge & base atom \\
beat & when & discrete synchronous time, collide then stream, invertible, no built-in direction & base atom \\
arrow & --- & the direction of time, the wake through the lean and the create move & emergent \\
attraction & --- & the radiation-pressure shadow of the active vacuum on the open growing mesh & emergent \\
\bottomrule
\end{tabular}
\caption{The ledger, eight base atoms and two emergent terms.}
\label{tab:ledger}
\end{table}
